% -------------------------------------------------------------------------
% WBS ID: RES-PHY-SYN
% Deliverable: Integrated Physical-Phenomena Model
% Status: Draft complete for regional 2D model assumptions
% Evidence status: Regional model physical assumptions established
% Review status: Not reviewed
% Last updated: 2026-07-07
% Dependencies: RES-PHY-SRC, RES-PHY-PRO, RES-PHY-NRS, RES-MOD-SPEC
% Downstream interfaces: RES-MOD-SPEC, RES-NUM-SPEC, RES-VER-SPEC
% -------------------------------------------------------------------------

\section{RES-PHY-SYN --- Integrated Physical-Phenomena Model}
\label{sec:res-phy-syn}

\subsection{Purpose and Scope}
\label{sec:res-phy-syn-purpose}

% TODO: Define what this Deliverable covers and explicitly excludes.

\subsection{Research Questions}
\label{sec:res-phy-syn-research-questions}

% TODO: State the questions that must be answered before the Deliverable
% can be accepted.

\subsection{Terminology and Definitions}
\label{sec:res-phy-syn-definitions}

% TODO: Define the technical terminology, symbols, quantities and
% classifications used in this Deliverable.

\subsection{Literature Evidence}
\label{sec:res-phy-syn-literature-evidence}

% TODO: Record evidence from peer-reviewed papers, authoritative datasets,
% standards, technical reports and primary documentation.
%
% Do not create a detached literature-summary list. Explain how each source
% supports, contradicts or limits a project decision.

\subsection{Governing Relations and Quantitative Definitions}
\label{sec:res-phy-syn-governing-relations}

% TODO: Record relevant equations, dimensionless groups, empirical models,
% extraction rules or quantitative definitions.
%
% Use align environments for displayed equations.
% Every displayed equation must have a unique \label{eq:...}.
% Do not place punctuation after displayed equations.

\subsection{Physical Mechanisms}
\label{sec:res-phy-syn-physical-mechanisms}

The regional 2D NLSWE model represents long-wave tsunami generation,
propagation, refraction over bathymetry, nearshore transformation,
run-up and inundation where depth-averaged hydrostatic assumptions are
acceptable for the selected outputs. The bed is a positive-upward scalar
field and may move dynamically during earthquake forcing.

The regional model does not resolve vertical velocity profiles,
nonhydrostatic impact pressure, air--water interface topology, local
breaking jets around a structure, structural deformation or FSI. Those
mechanisms belong to local 3D/interface/FSI workstreams.

\subsection{Characteristic Spatial and Temporal Scales}
\label{sec:res-phy-syn-characteristic-spatial-temporal-scales}

% TODO: Complete this RES-PHY Domain-specific subsection.

\subsection{Relevant Observables}
\label{sec:res-phy-syn-relevant-observables}

% TODO: Complete this RES-PHY Domain-specific subsection.

\subsection{Controlling Dimensionless Parameters}
\label{sec:res-phy-syn-controlling-dimensionless-parameters}

% TODO: Complete this RES-PHY Domain-specific subsection.

\subsection{Physical Regimes and Transition Conditions}
\label{sec:res-phy-syn-physical-regimes-transition-conditions}

% TODO: Complete this RES-PHY Domain-specific subsection.

\subsection{Implications for Model Validity}
\label{sec:res-phy-syn-model-validity-implications}

The regional mathematical model is established at research-specification
level, but physical credibility still depends on domain selection,
bathymetry/topography quality, source-data quality, wet/dry validation
and comparison against observations. Agreement in offshore propagation
does not by itself validate local structure loading.

\subsection{Candidate Approaches}
\label{sec:res-phy-syn-candidate-approaches}

% TODO: Define the credible candidate approaches considered.

\subsection{Critical Comparison}
\label{sec:res-phy-syn-critical-comparison}

% TODO: Compare candidates using physically and project-relevant criteria.
% Do not select an approach solely on popularity or implementation ease.

\subsection{Assumptions and Applicability}
\label{sec:res-phy-syn-assumptions}

% TODO: State assumptions and the conditions under which they are valid.

\subsection{Limitations and Failure Conditions}
\label{sec:res-phy-syn-limitations}

% TODO: State where the evidence, model or selected approach ceases to be
% reliable or applicable.

\subsection{Project-Specific Conclusions}
\label{sec:res-phy-syn-conclusions}

% TODO: Convert the evidence into conclusions specific to the Studentship
% project. Do not merely restate literature findings.

\subsection{Selected Approach and Rationale}
\label{sec:res-phy-syn-selected-approach}

% TODO: Record the selected approach only after sufficient evidence exists.
% State the alternatives rejected and the reasons for rejection.

The selected physical regional baseline is a depth-averaged,
hydrostatic, horizontal two-dimensional long-wave model. Optional
dispersive or Coriolis sensitivities may be assessed later if the
selected domain or validation evidence requires them.

\subsection{Unresolved Decisions}
\label{sec:res-phy-syn-unresolved-decisions}

% TODO: Record unresolved questions, required evidence and decision owners.

Open physical work is final case-study/domain selection, final source
and terrain datasets, validation of propagation and inundation, and
handoff conditions for local 3D/interface modelling.

\subsection{Requirements and Handoffs}
\label{sec:res-phy-syn-requirements}

% TODO: State requirements passed to other Research Domains, Software,
% verification activities or later execution work.

\subsection{Deliverable Completion Record}
\label{sec:res-phy-syn-completion-record}

% TODO: Record whether the Deliverable is:
%
% - Not started
% - In progress
% - Draft complete
% - Reviewed
% - Accepted
%
% Acceptance requires:
%
% - the research questions to be answered;
% - claims to be supported by appropriate evidence;
% - assumptions and limitations to be explicit;
% - the project-specific conclusion to be stated;
% - downstream requirements to be traceable;
% - unresolved decisions to be closed or formally deferred.
